% !TeX program = xelatex
\documentclass[letterpaper,10pt]{article}

\usepackage[margin=0.50in]{geometry}
\usepackage{iftex}
\usepackage{fontspec}
\setmainfont{Charter}
\usepackage[hidelinks]{hyperref}
\usepackage{enumitem}

% ATS-compatible text extraction under pdfLaTeX
\ifPDFTeX
  \input{glyphtounicode}
  \pdfgentounicode=1
\fi

\pagestyle{empty}
\setlength{\parindent}{0pt}
\setlength{\parskip}{0pt}
\urlstyle{same}
\raggedright

\setlist[itemize]{
  leftmargin=0.17in,
  itemsep=0.4pt,
  topsep=1.0pt,
  parsep=0pt,
  partopsep=0pt
}

\newcommand{\resumesection}[1]{%
  \vspace{3.1pt}%
  {\large\bfseries #1}\par
  \vspace{0.9pt}%
  \hrule height 0.4pt
  \vspace{2.0pt}%
}

\newcommand{\subheading}[4]{%
  \textbf{#1}\hfill #2\\
  \textit{#3}\hfill \textit{#4}\par
  \vspace{-1.5pt}%
}

\newcommand{\projectheading}[2]{%
  \textbf{#1}\hfill #2\par
  \vspace{-1.5pt}%
}

\begin{document}

% -------------------- Header --------------------
\begin{center}
  {\Huge\bfseries Emilio Daza}\\[-1pt]
  {\small
  Hanover, NH
  \;|\;
  +1 (603) 443-3235
  \;|\;
  \href{mailto:emilio.sebastian.daza.vigo.27@dartmouth.edu}{emilio.sebastian.daza.vigo.27@dartmouth.edu}
  \;|\;
  \href{https://emiliodaza.com}{emiliodaza.com}}
\end{center}
\vspace{-8pt}

{\small
Dartmouth computer science student graduating June 2027 with experience
building cross-platform applications, backend systems, intelligent
payment-routing models, and healthcare ML pipelines for heart-disease
prediction and brain-tumor MRI classification.
}

% -------------------- Education --------------------
\resumesection{Education}

\subheading
  {Dartmouth College}{Hanover, NH}
  {B.A. in Computer Science, Minor in Statistics}{Expected Jun. 2027}

{\small
\textbf{Relevant Coursework:}
Computer Architecture, Compilers, Discrete Mathematics, Robotics Design \&
Programming, Computer Vision, Multivariate Statistics \& Statistical Learning, and Probability \& Statistical Inference.\\
\textbf{Beca Cometa Scholar (2023):}
1 of 2 students selected nationwide in Peru for a full undergraduate scholarship
at a top U.S. university.
}

% -------------------- Technical Skills --------------------
\resumesection{Technical Skills}

{\small
\textbf{Languages:}
Python, TypeScript, JavaScript, SQL, Java, C, C++, Dart, R\\
\textbf{Frameworks:}
React Native, Expo, Flutter, React, PyTorch, Gymnasium, R-Studio\\
\textbf{Backend/Data:}
ClickHouse, PostgreSQL, Supabase, Firebase, REST APIs, Edge Functions,
row-level security\\
\textbf{Tools/Systems:}
AWS, Docker, Linux, Git, GitHub, NetworkX, ROS 2, Gazebo
}

% -------------------- Experience --------------------
\resumesection{Experience}

\subheading
  {Izipay}{Lima, Peru}
  {Machine Learning Engineer Intern}{Jun. 2026 -- Aug. 2026}

\begin{itemize}
  \item Developed and evaluated an RL-based intelligent payment-routing
  prototype in a synthetic simulator for Izipay, a leading Peruvian processor
  serving 1.2M+ merchants; now adapting the architecture by querying
  transaction-node data from an AWS-hosted ClickHouse instance using SQL.

  \item Modeled processor selection under outages, latency variation, technical
  failures, and payment rejection as a Markov decision process; developed a
  custom Gymnasium environment and PyTorch Double DQN agent over a NetworkX
  topology with action masking, experience replay, and target networks.

  \item Evaluated the learned policy across 1,000 seeded episodes generated
  from synthetic routing probabilities: 92.3\% simulated approval and
  178.4\,ms mean latency; automated JSON/Matplotlib reporting and presented
  results to senior strategy and business-solutions leaders.
\end{itemize}

\subheading
  {Algoverse AI Research Program}{Remote}
  {Machine Learning Researcher}{Jun. 2026 -- Sep. 2026}

\begin{itemize}
  \item Developing a formal research project on whether finite prefixes of
  ordinary RL training can distinguish reward misspecification from
  task-relevant observation insufficiency, and what minimum diagnostic
  interventions make the failure modes identifiable.

  \item Collaborating with a three-person ML research team under the mentorship
  of a Google Cloud AI \& Infrastructure Software Engineer; selected for a
  12-week publication-oriented program with a competitive 30\% tuition
  scholarship.
\end{itemize}

\subheading
  {Dartmouth Center for Technology and Behavioral Health}{Hanover, NH}
  {Software Engineer, Cross-Platform App}{Nov. 2025 -- Mar. 2026}

\begin{itemize}
  \item Implemented and validated Flutter desktop functionality on Arch Linux
  for Evergreen, a research-backed platform supporting personalized student
  well-being.

  \item Designed a dual-mode interface separating internal diagnostic and
  developer workflows from the participant-facing user experience.
\end{itemize}

\subheading
  {Dartmouth College, E.E. Just Research Program}{Hanover, NH}
  {Funded Machine Learning Research Intern}{Jun. 2024 -- Aug. 2024}

\begin{itemize}
  \item Selected as 1 of 8 first-year students for a funded summer program;
  built PyTorch feedforward and CNN pipelines for heart-disease prediction and
  four-class brain-tumor MRI classification, then presented the research at the largest math conference in the world: the Joint Mathematics Meeting in 2025.
\end{itemize}

% -------------------- Selected Projects --------------------
\resumesection{Selected Projects}

\projectheading
  {Thyme -- Co-Founder \& Founding Software Engineer}
  {Apr. 2026 -- Present}

\begin{itemize}
  \item Built a TypeScript React Native/Expo application for iOS and Android
  with design and public-health teammates; distributed an iOS beta through
  TestFlight and an Android APK for internal testing.

  \item Designed a Supabase/PostgreSQL backend with authentication, row-level
  security, Edge Functions, community and menu data, and remote version
  controls; implemented persistent profiles, settings, and nutrition logs.

  \item Integrated OpenAI APIs for multimodal meal and menu analysis, bilingual
  speech transcription and text-to-speech, and generated cooking guidance;
  developed location-aware dining discovery using Google Places and
  campus-menu data sources.
\end{itemize}

\projectheading
  {Robotics Mapping, Planning, and Perception Systems}
  {Spring 2026}

\begin{itemize}
  \item Implemented ROS 2 systems in Gazebo and Stage for odometry, coordinate
  transforms, LiDAR processing, and occupancy-grid mapping; a YOLO11n and FastSAM perception pipeline producing target bounding boxes and segmentation masks from camera video streams; BFS, DFS, and A* planners plus a LiDAR mapper
  using Bresenham raycasting, log-odds updates, and dynamically expanding
  occupancy grids.
\end{itemize}

\end{document}
